\scanpage{89}
Un homomorphisme $u:A\longrightarrow B$ de V.A. est connu quand on connaît la classe de cohomologie associée à son graphe \unsure{dans|de} $A\times B$,
\deleted{la \unsure{coho.}} (i.e. si on la connaît \unsure{en tant} \unsure{que} dans le \unsure{ann.} cohomologique $\ell$-adique), car alors on connaît le \unsure{graphe|groupe} $\Gamma\inserted{{}\subset A\times B}$ de $u$ car $T_{Q_\ell}(\Gamma)\subset T_{Q_\ell}(A\times B)$ est alors connu.

\uline{\unsure{D}} $=$, une ss-V.A. de $A$ est connue quand on connaît sa classe de cohomologie.

L'``interprétation'' cohomologique de l'isom
\[
C^1(A\times B)/\deleted{\unsure{=}}\ C^1(A)\times C^1(B)\simeq\operatorname{Hom}(A,B^*)
\]
est bien \unsure{clair|claire}, via
\diagram{\begin{tikzcd}[ampersand replacement=\&,column sep=large,row sep=large]
\bigwedge^2(H^1(A)\times H^1(B))(1)/\bigwedge^2 H^1(A)(1)\times\bigwedge^2 H^1(B)(1)
\arrow[r,phantom,"{\simeq}"] \& H^1(A)\otimes H^1(B)(1)\arrow[d,phantom,"{\unsure{\wr}}"]\\
\&\operatorname{Hom}(T_\ell(A),T_\ell(B^*))
\end{tikzcd}}

\uline{Question} Prouver que pour qu'une \unsure{classe} de $H^1(A)\otimes H^1(B)(1)\subset H^2(A\times B)(1)$ soit algébrique

\margin{88}
